\section{El funcional de una línea en el cuadrado marcado}\label{app:weights}
La segunda media de frontera compara la inserción conjunta con su aproximación al primer
orden. Insertar \(e\) cambia el peso de \(f\) porque elimina su propia fila y columna.
Ordenamos \(e,f\) por fila creciente y conservamos los dos indicadores cruzados \(X,Y\).
Aquí \(M_e\) es el menor simple de torres, sin la fila ni la columna de \(e\). Las primas
se evalúan en ese menor. Sean \(t=p'_f/p_f-1\), \(G_{ef}=\mathsf G(e,f)\) y
\(D=\hat c_e+\hat c_f-G_{ef}\). El producto exacto de los cocientes de inserción da
\[
 e^{u}=\frac{\pi^T_{ef}}{p_ep_f}
      =(1+\hat c_e)(1+\hat c'_f)(1+t)\exp(r_e+r'_f).
\]
La estimación de incrementos de torres \((A^+)\) del Apéndice~\ref{app:rook} es
\(r_i=-\mathsf E_i+O(j^{-3})\), también en el menor
simple. Con \(t=-G_{ef}+O(j^{-2})\) y \(\hat c'_f-\hat c_f=O(j^{-2})\)
(Corolario~\ref{cor:dos-casillas}), el resto físico dividido
por \(p_ep_f\) es \(\eta+O(j^{-3})\), donde
\begin{align*}
 \eta={}&(t+G_{ef})+(\hat c'_f-\hat c_f)+\hat c_e\hat c_f+t(\hat c_e+\hat c_f)\\
        &-\mathsf E_e-\mathsf E'_f+(p_e+p_f)G_{ef}\\
     ={}&B_{\rm phys}+D^2/2+O(j^{-3}),\\
 B_{\rm phys}={}&u-D+(p_e+p_f)G_{ef}.
\end{align*}
Estos siete sumandos valen en cada estado y orientación. El cuadrado \(D^2/2\) sólo requiere
el término lineal de la inserción. Para identificar el resto \(B_{\rm phys}\) hay que comparar la
acción y el prefactor en la máscara y en su menor doble. Para hacerlo en una misma familia
retiramos sólo \(e,f\) y conservamos sus cuatro líneas. Los dos márgenes de la casilla \(i\)
pasan a ser \(1-x_i\), y todos los demás valen uno. Sean \(F_-(x),\ell_-(x)\) la acción
estacionaria y la corrección de torres de esa máscara perforada. Ponemos
\(R(x)=F_-(x)+2\sum_i(1-x_i)\log(1-x_i)\) y \(A(x)=\ell_-(x)+\sum_i\log(1-x_i)\).
Ambas expresiones regularizadas llegan a \(x=(1,1)\): el sistema visible es el menor doble,
y la masa total de su Hessiana conserva la constante de tamaño correcta.
Sean \(p=(p_e,p_f)^\top\), \(\mathbf1=(1,1)^\top\) y
\(D_p=\operatorname{diag}(p_i(1-p_i))\).
Para la Hessiana perforada de márgenes \(T_-\) y la incidencia \(a_i\),
\(R''_{il}=2\delta_{il}/(1-x_i)-a_i^\top T_-^+a_l\).
Reinsertar las casillas añade \(\sum_i[p_i(1-p_i)/(1-2p_i)]a_i a_i^\top\).
El lema del determinante en el complemento del calibre y la estacionariedad dan
\begin{align*}
 u={}&R(\mathbf1)-R(p)-(\mathbf1-p)^\top\nabla R(p)\\
     &+A(\mathbf1)-A(p)+\tfrac12\log\det(I-R''(p)D_p).
\end{align*}
Sean \(\partial=\partial_e+\partial_f\), \(\mathcal H_{\rm par}=R''(0)\) y \(\Delta C=C(\mathbf1)-C(0)\).
La derivada de la ecuación estacionaria en los dos pesos da la respuesta
\((D_p^{-1}-R''(p))^{-1}\). Se desarrollan el determinante y las respuestas de
\(\hat c_e,\hat c_f,G_{ef}\). Los términos lineales en \(p_i (\mathcal H_{\rm par})_{il}\) se cancelan por
\(\mathbf1^\top \mathcal H_{\rm par}p=\sum_i p_i(\mathcal H_{\rm par})_{ii}+(p_e+p_f)(\mathcal H_{\rm
 par})_{ef}\). Queda
\[
 B_{\rm phys}=\Delta R-\partial R(0)-\partial^2 R(0)/2+\Delta A-\partial A(0)+O_b(j^{-3}).
\]
Para obtener los coeficientes, a cada línea marcada se le añaden dos copias de signos opuestos.
En la copia positiva \(\lambda_i=1-x_i\), y en la parte ordinaria y las copias negativas
\(\lambda_i=1\). Con masa auxiliar \(s\), la medida
de número es \(M=a+s\rho\) y la de margen \(N=\lambda M\), con \(\rho=0,-1,+1\) en la parte
ordinaria, las copias negativas y las positivas. En \(s=1/j\) las copias negativas eliminan
los originales exactamente. Para la incidencia binaria \(Y_{il}\) y sus factores de
Sinkhorn~\cite{sk}
\(W_{il}\), las ecuaciones normalizadas son
\[
 w_{il}(1-s\lambda_i\lambda_l w_{il})=W_{il}\exp(z_i+y_l),\qquad
 \int Y_{il}w_{il}\,dN_C(l)=\int Y_{il}w_{il}\,dN_R(i)=1 .
\]
Denotemos por \(R_q,A_q\) los coeficientes de \(s^q\) de la acción y del prefactor regularizados
sobre \(x_e=x_f=t\). En la acción, el coeficiente de \(s^0\) es afín y \(R_1\) tiene grado
a lo sumo dos en \(t\). En el prefactor \(A_0\) es constante. Los dos términos cuadráticos
del factor local y de la diagonal de márgenes se cancelan en cada extremo marcado, y \(A_1\)
es afín. Para controlar el resto a corte fijo, sea \(\centrado=YW-1\) el núcleo basal centrado de
la máscara perforada en la buena aproximación. La brecha \(\|\centrado\|_{2\to2}\le r_b<1\)
acota \((I-\centrado^*\centrado)^{-1}\). Las cotas cuadráticas de las secciones de \(\centrado\) acotan
 también
en norma de supremo la inversa del jacobiano de márgenes en el calibre. En las copias bastan sus
ecuaciones con diagonal uno. Para \(|x_i|\le2\), tanto \(\lambda\) como la variación total
de \(m=\lambda\rho\) están acotadas. La ecuación implícita da por ello una rama holomorfa
en un disco \(|s|<\varepsilon_b\) uniforme en los pares. En el mismo disco reducido, los
grados de la Hessiana y su masa siguen cerca de uno. La serie de Fredholm del determinante
centrado admite el mayorante de Hadamard. En el origen la brecha separa su determinante de
cero, pues la suma de cuadrados singulares es acotada. Si \(M_b\) domina el determinante
en el disco inicial y \(m_b>0\) minora su módulo en cero, Cauchy acota su derivada
en el disco mitad por \(2M_b/\varepsilon_b\). En el radio reducido
\(r_b'=\min(\varepsilon_b/2,m_b\varepsilon_b/(4M_b))\) se mantiene el módulo
\(\ge m_b/2\), lo que fija la rama basal del logaritmo. Los factores \(s^{-1},s^{-2}\)
requieren uno y dos coeficientes adicionales en sus numeradores. Para una función regularizada
acotada por \(K_b\) en \(|s|\le r<r_b'\), la cola desde grado tres, evaluada en \(1/j\),
es a lo sumo \(K_b(rj)^{-3}/(1-(rj)^{-1})\). Cauchy también en el parámetro
\(t\), en un disco interior a \(|t|<2\) que contenga \([0,1]\), da la misma tasa para
sus tres primeras derivadas. Todos estos restos son \(O_b(j^{-3})\).
Con primas para las derivadas en \(t\), se obtiene
\begin{align}
 L_s&=A'_1(0)+\tfrac12R''_1(0),\notag\\
 B_{2,s}&=\tfrac12\int_0^1(1-t)^2R'''_2(t)\,dt+\int_0^1(1-t)A''_2(t)\,dt .\label{eq:estado-def}
\end{align}
En \(L_s,B_{2,s}\) el subíndice designa el estado de los cruces. Los polinomios
\(R_1,R_2,A_1,A_2\) tienen grados a lo sumo \(4,6,6,8\), de modo que nueve nodos distintos
determinan sus derivadas algebraicamente. La evaluación en punto flotante requiere su propio
control de error. Damos las recurrencias en el perfil para cada orientación y estado.
El apéndice~\ref{app:transfer} justifica su uso al promediar la máscara física.

\subsection{Datos del par y recurrencia de la acción}\label{sec:par-accion}

Llamemos \(z_{ij}=(r+ih,c+jh)\), \(i,j\in\{0,1\}\), a las cuatro casillas del cuadrado
geométrico. Numeramos las filas raíz por \((r_0,r_1)=(r,r+h)\). En la orientación \(D\)
numeramos las columnas raíz por \((c_0,c_1)=(c,c+h)\) y en la orientación \(A\) por
\((c_0,c_1)=(c+h,c)\). Así se perforan siempre \((r_0,c_0),(r_1,c_1)\). En el orden
de raíces, el bloque de incidencias es \(\left(\begin{smallmatrix}0&X\\Y&0\end{smallmatrix}\right)\).
En coordenadas del tablero esto da la siguiente tabla.
\begin{center}
\begin{tabular}{@{}ccc@{}}
\toprule
Estado \((X,Y)\) & Cruces abiertos para \(D\) & Cruces abiertos para \(A\)\\
\midrule
\((0,0)\)& ninguno & ninguno\\
\((1,0)\)& \(z_{01}\) & \(z_{00}\)\\
\((0,1)\)& \(z_{10}\) & \(z_{11}\)\\
\((1,1)\)& \(z_{01},z_{10}\) & \(z_{00},z_{11}\)\\
\bottomrule
\end{tabular}
\end{center}
Cada fila de la tabla lleva el peso \(w_{o,b}q_{o,(X,Y)}\) del apéndice~\ref{app:boundary}.
Para construir la integral a corte \(b>0\), un tipo ordinario de fila es \((u,\xi)\),
donde \(u\in[0,1]\) y \(\xi=(\xi_0,\xi_1)\) indica las aristas que llegan a las dos columnas
raíz. Escribamos \(e_1(v)=v,\ e_0(v)=1-v\). Las medidas ordinarias de fila y columna son
\[
 d\alpha_R(u,\xi)=du\prod_{k=0}^1 e_{\xi_k}(v_b(u,c_k)),\qquad
 d\alpha_C(v,\zeta)=dv\prod_{k=0}^1 e_{\zeta_k}(v_b(r_k,v)).
\]
A cada lado se añaden dos copias de cada raíz, de masa basal cero. La medida firmada
\(\rho_\sigma\), en el lado \(\sigma=R,C\), asigna masas \(-1,+1\) a sus copias
negativa y positiva. La capacidad \(\lambda\) vale \(1-t\) en las copias positivas y uno
en los demás tipos. Fijamos \(m_\sigma=\lambda\rho_\sigma\),
\(M_\sigma=\alpha_\sigma+s\rho_\sigma\) y \(N_\sigma=\alpha_\sigma+s m_\sigma\).

Sea \(V_{il}\) la incidencia condicional entre un tipo fila \(i\) y un tipo columna \(l\).
Entre tipos ordinarios de posiciones \(u,v\) vale \(v_b(u,v)\). Entre una fila ordinaria y
una copia de la columna raíz \(k\) vale \(\xi_k\). Entre una copia de la fila raíz \(k\) y
una columna ordinaria vale \(\zeta_k\). Entre copias se usa la matriz del estado indicado.
En todos los bloques \(W=W_b^0=p_b/v_b\) se evalúa en las posiciones de ambos tipos.
Con \(q_{il}=\lambda_i\lambda_l,\ \kappa_{1,il}=q_{il}W_{il}\), sustituimos la incidencia por \(V\)
en las ecuaciones de márgenes de este apéndice. Para el determinante, la corrección por
la ida y vuelta en una misma arista se calcula en la sección~\ref{sec:par-prefactor}.

\paragraph{Potenciales y densidades por orden.}
Sean \(\mathcal Z_d=(Z_d,H_d)=[s^d](z,y)\) y \(h_d(i,l)=Z_d(i)+H_d(l)\).
Partimos de \(\mathcal Z_0=0,\ w_0=W\). Para \(d=1,2,3\), el término conocido \(v_d\) de
\(w_d=[s^d]w=Wh_d+v_d\) se obtiene de \(w(1-sqw)=W\exp(z+y)\).
Omitiendo los índices de arista,
\begin{align*}
 v_1/W&=\kappa_1,\\
 v_2/W&=h_1^2/2+2\kappa_1 h_1+2\kappa_1^2,\\
 v_3/W&=h_1h_2+h_1^3/6+\kappa_1(2h_2+2h_1^2)+6\kappa_1^2h_1+5\kappa_1^3 .
\end{align*}
A orden \(d\), los dos márgenes y el calibre determinan \(\mathcal Z_d\) mediante el sistema
\begin{align*}
 Z_d(i)+\int VW H_d(l)\,d\alpha_C(l)
   &=-\int Vv_d\,d\alpha_C-\int Vw_{d-1}\,dm_C ,\\
 H_d(l)+\int VW Z_d(i)\,d\alpha_R(i)
   &=-\int Vv_d\,d\alpha_R-\int Vw_{d-1}\,dm_R ,\\
 \int Z_d\,d\alpha_R&=\int H_d\,d\alpha_C .
\end{align*}
Se resuelve primero el bloque ordinario. La brecha de \(p_b\) da la inversa en el calibre,
y las incógnitas de las copias se recuperan de su ecuación con diagonal uno.
Los lados derechos son compatibles porque las sumas de los dos márgenes respecto
de \(N_R,N_C\) coinciden y estas medidas tienen la misma masa.

\paragraph{Extracción de la acción.}
Escribimos \(\langle\alpha,\mathcal Z_d\rangle=\int Z_d\,d\alpha_R+\int H_d\,d\alpha_C\)
y análogamente para \(m\). Sean \(D_j=[s^j](N_R\otimes N_C)\). En particular,
\(D_0=\alpha_R\otimes\alpha_C,\ D_1=m_R\otimes\alpha_C+\alpha_R\otimes m_C\) y
\(D_2=m_R\otimes m_C\). En cada arista se calculan
\[
 f_1=\tfrac12 qW^2,\qquad f_2=qWw_1+\tfrac13q^2W^3,\qquad
 f_3=q(Ww_2+w_1^2/2)+q^2W^2w_1+q^3W^4/4 .
\]
En la acción dual, extraer los términos de grados uno y dos da
\begin{align*}
 R_1(t)&=-\langle\alpha,\mathcal Z_2\rangle-\langle m,\mathcal Z_1\rangle
          -\int Vf_2\,dD_0-\int Vf_1\,dD_1,\\
 R_2(t)&=-\langle\alpha,\mathcal Z_3\rangle-\langle m,\mathcal Z_2\rangle
          -\int Vf_3\,dD_0-\int Vf_2\,dD_1-\int Vf_1\,dD_2 .
\end{align*}
En efecto, \(f_1,f_2,f_3\) son los tres primeros coeficientes de
\(\sum_{a\ge2}s^{a-2}q^{a-1}w^a/a\), que se integra con signo negativo contra
\(V\,d(N_R\otimes N_C)\). La contribución de los potenciales es
\(-s^{-1}(\int z\,dN_R+\int y\,dN_C)\), módulo los términos afines en \(t\)
que se anulan en el resto de Taylor.

\subsection{Recurrencia del prefactor y traza de una arista}\label{sec:par-prefactor}

Con los \(w_d\) de la sección~\ref{sec:par-accion}, sean \(p=sqw\) y
\(\ell=w(1-p)/(1-2p)\). El grado de una fila es
\(d_R(i)=\int V_{il}\ell_{il}\,dN_C(l)\) y el de una columna se define transponiendo
la fórmula. Ambos valen uno en \(s=0\). La masa total
\(g=\int d_R\,dN_R=\int d_C\,dN_C\) también empieza en uno. Para calcular
\(d_\sigma=1+\sum_{a\ge1}s^a d_{\sigma,a}\), desarrollamos
\(\ell-w=sqw^2/(1-2sqw)\). Con \(\kappa_1=qW\), sus tres primeros coeficientes son
\begin{align*}
 E_1&=W\kappa_1,\qquad E_2=2\kappa_1 w_1+2W\kappa_1^2,\\
 E_3&=q(2Ww_2+w_1^2)+6\kappa_1^2w_1+4W\kappa_1^3 .
\end{align*}
Así \(\ell_0=W,\ \ell_1=w_1+E_1,\ \ell_2=w_2+E_2\). Con \(E_0=0\), los márgenes
\(\int Vw\,dN_{\rm opuesto}=1\) dan para \(a=1,2,3\)
\[
 d_{R,a}=\int VE_a\,d\alpha_C+\int VE_{a-1}\,dm_C,\qquad
 d_{C,a}=\int VE_a\,d\alpha_R+\int VE_{a-1}\,dm_R .
\]
Ponemos \(g_a=\int d_{R,a}\,d\alpha_R+\int d_{R,a-1}\,dm_R\) para \(a=1,2\),
con \(d_{R,0}=1\). Para cada lado, los coeficientes \(D_{\sigma,a}^{\log}=[s^a]\log d_\sigma\)
son \(d_1,\ d_2-d_1^2/2,\ d_3-d_1d_2+d_1^3/3\) en órdenes uno, dos y tres,
respectivamente, con \(d_a=d_{\sigma,a}\).

\paragraph{Términos locales y determinante.}
Escribamos \(\mu_a=[s^a](M_R\otimes M_C)\), con \(\mu_a=0\) para \(a\notin\{0,1,2\}\).
Para \(\kappa_1=qW,\ \kappa_2=qw_1,\ \kappa_3=qw_2\), los coeficientes de
\(-\log(1-2p)+2\log(1-p)\) en órdenes dos, tres y cuatro son
\[
 C_2=\kappa_1^2,\qquad C_3=2\kappa_1\kappa_2+2\kappa_1^3,\qquad
 C_4=\kappa_2^2+2\kappa_1\kappa_3+6\kappa_1^2\kappa_2+\tfrac72\kappa_1^4 .
\]
Para \(a=1,2\), las aristas, la masa total y las diagonales de la Hessiana aportan
\begin{align*}
 H_a^{\rm loc}={}&\tfrac12\sum_{j=2}^{a+2}\int V C_j\,d\mu_{a+2-j}
 +\tfrac12[s^a]\log(1+sg_1+s^2g_2)\\
 &-\tfrac12\sum_{\sigma=R,C}
 \left(\int D_{\sigma,a+1}^{\log}\,d\alpha_\sigma+
       \int D_{\sigma,a}^{\log}\,d\rho_\sigma\right).
\end{align*}
Las distintas medidas proceden de la masa de la Hessiana \(g\), integrada contra \(N_R\),
y de la suma de sus logaritmos de grado, integrada contra la medida de número \(M_\sigma\).

Para los bloques centrados, los operadores \(K_a,T_a\), \(a=0,1,2\), se obtienen
extrayendo el coeficiente de cada núcleo junto con su medida de integración.
\begin{align*}
 K_a(i,dl)&=[s^a]\left\{\left(\frac{V_{il}\ell_{il}}{d_R(i)}-\frac{d_C(l)}g\right)
                       (d\alpha_C(l)+s\,dm_C(l))\right\},\\
 T_a(l,di)&=[s^a]\left\{\left(\frac{V_{il}\ell_{il}}{d_C(l)}-\frac{d_R(i)}g\right)
                       (d\alpha_R(i)+s\,dm_R(i))\right\}.
\end{align*}
Para evaluar los inversos escalares
se usa \((1+sa+s^2b)^{-1}=1-sa+s^2(a^2-b)+O(s^3)\). Sean
\[
 I_0=(I-K_0T_0)^{-1},\qquad X_1=K_1T_0+K_0T_1,\qquad
 X_2=K_2T_0+K_1T_1+K_0T_2 .
\]
El desarrollo de \(-\tfrac12\log\det(I-KT)\) para los dos bloques de medias da
\[
 H_1^{\rm tr}=\tfrac12\operatorname{tr}(I_0X_1),\qquad
 H_2^{\rm tr}=\tfrac12\operatorname{tr}(I_0X_2)+\tfrac14\operatorname{tr}((I_0X_1)^2).
\]

\paragraph{La ida y la vuelta por una casilla ordinaria.}
En una traza de longitud dos figura el mismo indicador \(I\) en ambas direcciones.
Para una arista entre tipos ordinarios, \(\mathbb E[I^2]=V\) y
\(\mathbb E[I]^2=V^2\), condicionalmente a los tipos. Al multiplicar los dos núcleos
centrados, los términos constantes y lineales tienen las medias correctas. En el
término cuadrático falta \(V(1-V)\ell^2/(d_Rd_C)\). Sea \(d\mu_*=\tfrac12
V(1-V)\,d\alpha_R\,d\alpha_C\), incluido el peso \(1/2\) de la traza. Para
\(r=d_{R,1}, c=d_{C,1}, r_2=d_{R,2}, c_2=d_{C,2}\), las dos contribuciones son
\begin{align*}
 H_1^{\rm rep}&=\int \big[2W\ell_1-(r+c)W^2\big]\,d\mu_*,\\
 H_2^{\rm rep}&=\int \big[\ell_1^2+2W\ell_2-2(r+c)W\ell_1
              +(r^2+c^2+rc-r_2-c_2)W^2\big]\,d\mu_* .
\end{align*}
Los tipos raíz tienen incidencia fijada. Todas sus repeticiones usan ese mismo bit.
Concluimos que, para cada orientación y estado, el prefactor pedido es
\[
 A_a(t)=H_a^{\rm loc}+H_a^{\rm tr}+H_a^{\rm rep},\qquad a=1,2 .
\]

\paragraph{Ejemplo de un sumando de \(A_2(1)\).}
Tomemos como dato de prueba el núcleo constante \(v=1/2\), cuyo escalado es \(p_b=1\).
Fijemos el estado \(X=Y=0\) y la capacidad \(1-t=0\). En ambos lados los patrones
ordinarios \(00,10,01,11\) tienen masa \(1/4\), y \(W=2\). Consideremos la ida y vuelta
entre una fila de patrón \(00\) y una columna también de patrón \(00\).
Para comprobar las entradas de la recurrencia, por simetría \(Z_a=H_a\). En el orden de
los cuatro patrones seguido de las dos copias negativas y las dos positivas, los sistemas dan
\[
 Z_1=(-2,0,0,2,-3,-3,-1,-1),\qquad
 Z_2=(2,-6,-6,-10,11/2,11/2,15/2,15/2).
\]
En la arista elegida resultan \(w_1=-4,\ w_2=8,\ \ell_1=0,\ \ell_2=8\).
En ambos extremos \(d_1=2,\ d_2=8\). De aquí
\[
 [s^2]\frac{\ell^2}{d_Rd_C}
 =[s^2]\frac{(2+8s^2)^2}{(1+2s+8s^2)^2}=16 .
\]
El sumando de \(H_2^{\rm rep}\) para estos dos patrones es, pues,
\(\tfrac12(1/4)^2(1/2)(1-1/2)\,16=1/8\).
La integración en las dos posiciones tiene masa uno porque el dato de prueba es constante.

